\begin{thebibliography}{99}

\bibitem{BRY2018} M.~Baake, J.~A.~G.~Roberts and R.~Yassawi, Reversing and extended symmetries of shift spaces, \emph{Discrete Contin. Dyn. Syst.} \textbf{38} (2018), 835--866.

\bibitem{CaDiCaL} A.~Biere, K.~Fazekas, M.~Fleury and M.~Heisinger, CaDiCaL, Kissat, Paracooba, Plingeling and Treengeling entering the SAT Competition 2020, in \emph{Proc. SAT Competition 2020}, Univ. Helsinki, 2020, 51--53.

\bibitem{BT2000} F.~Blanchard and P.~Tisseur, Some properties of cellular automata with equicontinuity points, \emph{Ann. Inst. H. Poincar\'e Probab. Statist.} \textbf{36} (2000), 569--582.

\bibitem{BNR1991} N.~Boccara, J.~Nasser and M.~Roger, Particlelike structures and their interactions in spatiotemporal patterns generated by one-dimensional deterministic cellular-automaton rules, \emph{Phys. Rev. A} \textbf{44} (1991), 866--875.

\bibitem{BLR1988} M.~Boyle, D.~Lind and D.~Rudolph, The automorphism group of a shift of finite type, \emph{Trans. Amer. Math. Soc.} \textbf{306} (1988), 71--114.

\bibitem{CFK2019} V.~Cyr, J.~Franks and B.~Kra, The spacetime of a shift endomorphism, \emph{Trans. Amer. Math. Soc.} \textbf{371} (2019), 461--488.

\bibitem{Hedlund1969} G.~A.~Hedlund, Endomorphisms and automorphisms of the shift dynamical system, \emph{Math. Systems Theory} \textbf{3} (1969), 320--375.

\bibitem{Hochman2010} M.~Hochman, On the automorphism groups of multidimensional shifts of finite type, \emph{Ergodic Theory Dynam. Systems} \textbf{30} (2010), 809--840.

\bibitem{PySAT} A.~Ignatiev, A.~Morgado and J.~Marques-Silva, PySAT: a Python toolkit for prototyping with SAT oracles, in \emph{Theory and Applications of Satisfiability Testing (SAT 2018)}, Lecture Notes in Comput. Sci. 10929, Springer, 2018, 428--437.

\bibitem{KS1989} B.~Kitchens and K.~Schmidt, Automorphisms of compact groups, \emph{Ergodic Theory Dynam. Systems} \textbf{9} (1989), 691--735.

\bibitem{KS1992} B.~Kitchens and K.~Schmidt, Markov subgroups of $(\Z/2\Z)^{\Z^2}$, in \emph{Symbolic Dynamics and its Applications}, Contemp. Math. 135, Amer. Math. Soc., 1992, 265--283.

\bibitem{KS2000} B.~Kitchens and K.~Schmidt, Isomorphism rigidity of irreducible algebraic $\mathbb{Z}^d$-actions, \emph{Invent. Math.} \textbf{142} (2000), 559--577.

\bibitem{Kopra2022} J.~Kopra, Glider automata on all transitive sofic shifts, \emph{Ergodic Theory Dynam. Systems} \textbf{42} (2022), 3716--3744.

\bibitem{Kurka1997} P.~K\r{u}rka, Languages, equicontinuity and attractors in cellular automata, \emph{Ergodic Theory Dynam. Systems} \textbf{17} (1997), 417--433.

\bibitem{Ledrappier1978} F.~Ledrappier, Un champ markovien peut \^etre d'entropie nulle et m\'elangeant, \emph{C. R. Acad. Sci. Paris S\'er. A-B} \textbf{287} (1978), A561--A563.

\bibitem{Ryan1972} J.~P.~Ryan, The shift and commutativity, \emph{Math. Systems Theory} \textbf{6} (1972), 82--85.

\bibitem{Salo2020} V.~Salo, Conjugacy of reversible cellular automata and one-head machines, preprint, arXiv:2011.07827 (2020).

\bibitem{Schmidt2002} K.~Schmidt, Algebraic $\mathbb{Z}^d$-actions, PIMS Distinguished Chair Lectures, University of Victoria, 2002.

\bibitem{Silverman2009} J.~H.~Silverman, \emph{The Arithmetic of Elliptic Curves}, 2nd ed., Grad. Texts in Math. 106, Springer, New York, 2009.

\bibitem{Takesue1987} S.~Takesue, Reversible cellular automata and statistical mechanics, \emph{Phys. Rev. Lett.} \textbf{59} (1987), 2499--2502.

\bibitem{Takesue1989} S.~Takesue, Ergodic properties and thermodynamic behavior of elementary reversible cellular automata. I. Basic properties, \emph{J. Stat. Phys.} \textbf{56} (1989), 371--402.

\bibitem{TM1990} T.~Toffoli and N.~Margolus, Invertible cellular automata: a review, \emph{Phys. D} \textbf{45} (1990), 229--253.

\bibitem{Ward1994} T.~Ward, Automorphisms of $\mathbb{Z}^d$-subshifts of finite type, \emph{Indag. Math. (N.S.)} \textbf{5} (1994), 495--504.

\bibitem{Salo2019Ryan} V.~Salo, Transitive action on finite points of a full shift and a finitary Ryan's theorem, \emph{Ergodic Theory Dynam. Systems} \textbf{39} (2019), 1637--1667.

\bibitem{KimRoush1990} K.~H.~Kim and F.~W.~Roush, On the automorphism groups of subshifts, \emph{Pure Math. Appl. Ser. B} \textbf{1} (1990), no.~4, 203--230.

\bibitem{MAM2013} G.~J.~Mart\'inez, A.~Adamatzky, F.~Chen and L.~Chua, On soliton collisions between localizations in complex elementary cellular automata: rules 54 and 110 and beyond, \emph{Complex Systems} \textbf{21} (2012), 117--142; arXiv:1301.6258.

\bibitem{BBGP1993} A.~Bobenko, M.~Bordemann, C.~Gunn and U.~Pinkall, On two integrable cellular automata, \emph{Comm. Math. Phys.} \textbf{158} (1993), 127--134.

\bibitem{BKP2021} B.~Bu\v{c}a, K.~Klobas and T.~Prosen, Rule 54: exactly solvable model of nonequilibrium statistical mechanics, \emph{J. Stat. Mech.} (2021), 074001.

\bibitem{Wolfram1984} S.~Wolfram, Universality and complexity in cellular automata, \emph{Phys. D} \textbf{10} (1984), 1--35.

\bibitem{Wolfram2002} S.~Wolfram, \emph{A New Kind of Science}, Wolfram Media, Champaign, IL, 2002.

\bibitem{PaperA} M.~Dayah, Reversible rule 110 has no hidden symmetries: six families of second-order cellular automata settled by edge moves, in preparation.

\bibitem{PaperB} M.~Dayah, Fronts at the speed of light: rigidity of permutive second-order cellular automata, in preparation.

\bibitem{PaperD} M.~Dayah, The 88 second-order elementary cellular automata: a classification by automorphism group, in preparation.

\bibitem{PaperE} M.~Dayah, The cellular automata that commute with Rule 54, and rigidity of the second-order automaton 94R, in preparation.

\end{thebibliography}
